\documentclass[10pt,a4paper]{article}
\usepackage[utf8]{inputenc}
\usepackage[T1]{fontenc}
\usepackage[french]{babel}
\usepackage{lmodern}
\usepackage[margin=1.9cm]{geometry}
\usepackage{amsmath,amssymb}
\usepackage{booktabs,array,tabularx,longtable,ltablex}
\keepXColumns
\usepackage{xcolor}
\usepackage{fancyvrb}
\usepackage{fvextra}
\usepackage{tikz}
\usetikzlibrary{arrows.meta,positioning}
\usepackage[hidelinks]{hyperref}
\usepackage{textcomp}
\usepackage{siunitx}
\usepackage{newunicodechar}
\newunicodechar{−}{-}
\newunicodechar{µ}{\textmu}
\newunicodechar{₂}{\textsubscript{2}}
\newunicodechar{→}{->}
\newunicodechar{≈}{~}
\sisetup{locale=FR,detect-all}

\setlength{\parskip}{0.4em}
\setlength{\parindent}{0pt}
\fvset{fontsize=\small,frame=single,framesep=2mm,rulecolor=\color{black!30}}
\newcommand{\vp}[1]{\text{\texttt{\detokenize{#1}}}}
\newcommand{\evC}{\textcolor{orange!80!black}{\textbf{C}}}
\newcommand{\evD}{\textcolor{red!70!black}{\textbf{D}}}
\newcolumntype{L}{>{\raggedright\arraybackslash}X}
\newcolumntype{P}[1]{>{\raggedright\arraybackslash}p{#1}}
\renewcommand{\arraystretch}{1.15}
\newcommand{\dd}{\mathrm{d}}

\title{Module d'anesthésie en temps réel\\[0.3em]
\large De la prémédication au réveil, avec la douleur comme variable d'état}
\author{Sébastien Lefebvre}
\date{19 novembre 2026 (version 1.0 du module)}

\begin{document}
\hypersetup{pdfauthor={Sébastien Lefebvre}}
\maketitle

\begin{abstract}
Ce document décrit le module d'anesthésie générale canine du moteur \emph{vetsim} tel qu'il est implémenté :
architecture, équations, modules JSON, provenance des valeurs, calibration, validation, API et interface web.
Il fait suite à la proposition du même jour (\vp{docs/PROPOSITION_ANESTHESIE.tex}), dont il réalise les jalons A0 à A5 et une partie de A6.
Les valeurs publiées sont recensées dans \vp{docs/ANESTHESIE_SOURCES.md}. Beaucoup de constantes restent des placeholders \evD{} :
le module sert à enseigner des mécanismes, jamais à décider d'une anesthésie réelle.
\end{abstract}

\tableofcontents

\section{Architecture}

Le moteur cardiovasculaire et métabolique (\vp{engine::evaluate}) est quasi stationnaire et reste la référence normative ;
il n'a pas été modifié. Une \vp{Session} (\vp{crates/vetsim-core/src/anesthesia/mod.rs}) porte les états dynamiques et avance par pas
$\Delta t \le \SI{60}{s}$. À chaque pas :
\begin{enumerate}
\item les états donnent des paramètres modifiés, toujours par les effets JSON du DSL (aucune formule clinique codée en dur) ;
\item l'évaluation existante calcule baroréflexe, hémodynamique, gaz du sang et décompensation (décès compris) ;
\item le débit cardiaque et la ventilation alvéolaire de cette évaluation font avancer les états jusqu'à $t + \Delta t$
(Runge-Kutta d'ordre 4 à pas fixe de \SI{0,5}{s} au plus, déterministe).
\end{enumerate}

\begin{center}
\begin{tikzpicture}[node distance=6mm and 8mm, every node/.style={font=\small},
  b/.style={draw, rounded corners, align=center, minimum height=9mm, inner sep=3pt},
  a/.style={-{Latex[length=2mm]}, thick}]
\node[b, fill=blue!6] (ev) {Actions\\(bolus, perfusion, vaporisateur,\\ventilation, chirurgie, fluides)};
\node[b, right=of ev, fill=orange!8] (st) {États\\PK + $C_e$, halogénés,\\CO$_2$/O$_2$, $\nu$, $S$, T°, remplissage};
\node[b, right=of st, fill=green!7] (pa) {Paramètres\\(effets JSON)};
\node[b, below=12mm of pa, fill=gray!10] (en) {\vp{evaluate}\\(inchangé)};
\node[b, below=of en, fill=red!6] (mo) {Instantané : moniteur,\\réflexes, alarmes, feuille};
\draw[a] (ev) -- (st); \draw[a] (st) -- (pa); \draw[a] (pa) -- (en); \draw[a] (en) -- (mo);
\draw[a] (en.west) -- node[below, font=\scriptsize]{DC, $\dot V_A$ (états à $t+\Delta t$)} (en.west -| st.south) -- (st.south);
\end{tikzpicture}
\end{center}

\paragraph{Invariant d'équilibre.} Une session sans médicament reproduit \vp{eval} : FC, PAM, PaCO$_2$ et PaO$_2$ identiques
à $10^{-9}$ près en relatif (test \vp{session_without_drugs_reproduces_eval}). La couche dynamique ne modifie donc pas le modèle validé.

\section{Équations}

\subsection{Pharmacocinétique}
Par agent : dépôts IM et SC, compartiments central et périphérique (constantes de \vp{pkpd::pk_constants}, donc allométrie et
fonctions rénale et hépatique du patient), site effet :
\begin{align*}
\dot A_{IM} &= -k_{a,IM} A_{IM}, \qquad \dot A_{SC} = -k_{a,SC} A_{SC},\\
\dot A_c &= k_{a,IM}A_{IM} + k_{a,SC}A_{SC} + R_{CRI} - \tfrac{CL + Q}{V_c}A_c + \tfrac{Q}{V_p}A_p, \qquad
\dot A_p = \tfrac{Q}{V_c}A_c - \tfrac{Q}{V_p}A_p,\\
\dot C_e &= k_{e0}\,(A_c/V_c - C_e).
\end{align*}
Un bolus IV est injecté à débit constant pendant sa durée (\SI{30}{s} par défaut, \vp{iv_push_s}, réglable par bolus) :
une injection rapide crée un pic plasmatique plus haut, donc plus d'apnée et d'hypotension (Murison 2001, Bigby 2017).
Les effets JSON du médicament reçoivent $E$ (Emax au site effet) et $EP$ (Emax plasmatique) ; les termes toxiques reçoivent $T$ (plasma).
Les effets hémodynamiques du propofol et de l'alfaxalone sont plasmatiques ($EP$) : le nadir de PAM survient à 55 s, avant l'équilibre du site effet
(intubation à 140 s, Cattai 2018).
Un antagoniste multiplie toutes les EC50 de ses cibles par $1 + \sum C_{e,ant}/K_i$ (antagonisme compétitif).

\subsection{Stocks de CO$_2$ et d'O$_2$}
\[
C_{CO_2}\,\frac{\dd P_aCO_2}{\dd t} = \dot V_{CO_2} - \frac{\dot V_A\,P_aCO_2}{0{,}863},\qquad
V_{CRF}\,\frac{\dd P_AO_2}{\dd t} = \dot V_A\,(P_IO_2 - P_AO_2) - 0{,}863\,\frac{\dot V_{CO_2}}{QR}.
\]
Leurs équilibres sont exactement les relations algébriques de \vp{gas.rs} ($P_aCO_2 = 0{,}863\,\dot V_{CO_2}/\dot V_A$,
$P_AO_2 = F_IO_2(P_B - 47) - P_aCO_2/QR$). L'évaluation reçoit une commande ventilatoire et une FiO$_2$ \emph{équivalentes},
qui lui font retrouver les valeurs dynamiques :
\[
\vp{vent_drive}_{eq} = \frac{0{,}863\,\dot V_{CO_2}/P_aCO_2}{\vp{va0}\cdot bw\cdot f_{K}\cdot \chi(P_aCO_2, pH, P_aO_2)},\qquad
F_IO_{2,eq} = \frac{P_AO_2 + P_aCO_2/QR}{P_B - 47},
\]
où $\chi$ est la stimulation chimique de \vp{gas_exchange} et $f_K$ la faiblesse hypokaliémique de \vp{evaluate}.
Le propofol et l'alfaxalone abolissent la commande ventilatoire au pic plasmatique ($\max(0,\,1 - a\,EP)$) : une injection rapide
provoque une apnée que le CO$_2$ ne lève qu'une fois la concentration redescendue, une injection lente la conserve (Bigby 2017, Murison 2001).
Après une hyperventilation contrôlée, la reprise spontanée attend que la PaCO$_2$ franchisse le seuil apnéique (stimulation chimique).
Dette d'O$_2$ : elle s'accroît quand l'apport passe sous l'apport critique (\vp{do2_crit_ml_kg_min}, consommation dépendante de
l'apport) ou sous la consommation, $\dot D = \max(VO_2 - DO_2,\ VO_2(1 - DO_2/DO_{2,crit}))^+$, et se rembourse sinon
(\vp{o2_debt_tau_min}). Sa fraction de la dette tolérée (\vp{o2_debt_tol_ml_kg}) déprime la contractilité et la FC
(effets JSON \vp{hypoxia_effects}) ; au-delà, arrêt hypoxique. Alarme « hypoxie tissulaire ». Placeholders \evD.
La réserve d'O$_2$ comprend la CRF et une fraction du sang (\vp{blood_o2_store_frac}) dont le contenu suit l'alvéole :
$(V_{CRF} + 0{,}863\,f_b\,V_S\,\mathrm{d}C_aO_2/\mathrm{d}P)\,\dot P_AO_2 = \dot V_A(P_IO_2 - P_AO_2) - 0{,}863\,\dot V_{CO_2}/QR$,
sans effet à l'équilibre. Apnée totale sous O$_2$ pur (isoflurane 1,3 MAC, après 10 min de ventilation contrôlée ; atélectasie comprise depuis le 20/11/2026) : SaO$_2$ < 90\,\% à 4,8 min, arrêt hypoxique à 10,3 min
(DeBehnke 1995 : 11,4 ± 2,4 min, PaCO$_2$ 93) ; préoxygénation au masque publiée : 187 à 298 s selon la PaO$_2$ de départ
(Ambros 2018, McNally 2009). Flow-by : 88 s (FiO$_2$ 0,40 du modèle ; 66 s à FiO$_2$ 0,30 chez Ambros). À l'air en apnée stricte : 28 s pour la SaO$_2$ artérielle (McNally 2009 : PaO$_2$ 62 mmHg et SaO$_2$ 82\,\% à 30 s ; la SpO$_2$ ne passe sous 90\,\% qu'à 70 s, retard du capteur absent du modèle).
Atélectasie : sous anesthésie, le gradient A-a reçoit $\vp{atelectasis_aa_mmhg}\cdot a\cdot\max(0, (F_IO_{2,\mathrm{eq}} - 0{,}21)/0{,}79)$ (110 mmHg, calé sur Itami 2019 et Kato 2024 : PaO$_2$ 523–550 mmHg sous O$_2$), avec la FiO$_2$ équivalente à la PAO$_2$ dynamique (24/11/2026) : égale à la FiO$_2$ réglée à l'état stable, elle évite qu'une intubation sous O$_2$ retranche le gradient d'une alvéole à 600 mmHg alors que la PAO$_2$ est encore basse. La ventilation réelle vaut $\vp{va0}\cdot bw\cdot\vp{vent_drive}\cdot f_K\cdot\chi$ en ventilation spontanée
et $\mathrm{FR}\,(V_T - V_D)$ en ventilation contrôlée. Apnée signalée si $\dot V_A < 10\,\%$ du repos.

\subsection{Halogénés}
Pressions partielles en \% d'atmosphère ; circuit filtre de volume $V_{circ}$, alvéoles ($V_{CRF}$), trois tissus perfusés
(groupe richement vascularisé dont le cerveau, muscle, graisse) :
\begin{align*}
V_{circ}\,\dot P_I &= \dot V_{DGF}(P_{vap} - P_I) - \dot V_A(P_I - P_A) \quad(\text{intubé}),\\
V_{CRF}\,\dot P_A &= \dot V_A(P_I - P_A) - \lambda_{s/g}\,\mathrm{DC}\,(P_A - P_{\bar v}), \qquad P_{\bar v} = \textstyle\sum_i f_i P_i,\\
\dot P_i &= \frac{f_i\,\mathrm{DC}}{V_i\,\lambda_{t/s,i}}\,(P_A - P_i).
\end{align*}
En circuit sans réinhalation, $P_I$ suit $P_{vap}$ (constante \SI{0,1}{min}). Non intubé, $P_I = 0$ (élimination à l'air).
Les effets de l'agent sont des effets JSON des variables $M = P_A/\mathrm{MAC}$ et $MB = P_{cerveau}/\mathrm{MAC}$ (contractilité, vasodilatation,
émoussement du baroréflexe, commande ventilatoire) ; le couplage avec
DC et $\dot V_A$ fait émerger la spirale du surdosage, le réveil retardé par l'hypoventilation et l'approfondissement lent en bas débit.

\subsection{Profondeur et conscience}
\[
H = \sum_{\text{halogénés}} \frac{P_{cerveau}}{\mathrm{MAC}} + \sum_{\text{IV}} \frac{C_e}{c_{imm}},\qquad
\Delta = \frac{H}{(1 - R)\,f_T},\qquad
L = \frac{H}{(1 - S_d)\,f_T},
\]
$R = 1 - \prod(1 - R_j)$ l'épargne de MAC combinée, $S_d = 1 - \prod(1 - S_j)$ la sédation combinée,
$f_T = \max(0{,}3;\ 1 - \vp{mac_per_c}\,(T_0 - T))$. Conscience perdue si $L \ge \vp{loc_mac}$, $L$ étant calculé avec un indice hypnotique où chaque halogéné pèse $\vp{loc_mac}/r_\mathrm{agent}$, $r_\mathrm{agent}$ = MAC-éveil/MAC (\vp{mac_awake_ratio} : isoflurane 0,77, sévoflurane 0,62, Hofmeister 2008 ; 24/11/2026) ; niveau de conscience
$(1 - S_d)\,(1 - x^n)$ avec $x = \mathrm{clip}(L/\vp{loc_mac}, 0, 1)$ et $n$ = \vp{awareness_hill} = 3 (calé sur les délais de réveil, Rastas 2025, Lozano 2009 ; voir \vp{VALIDATION_HORS_CALIBRAGE.tex}). L'intubation est refusée tant que la déglutition est présente ($L < \vp{swallow_mac}$). Plans et réflexes par seuils de $\Delta$ (palpébral, mâchoire, globe, mydriase)
et de $L$ (déglutition) ; un agent dissociatif majoritaire conserve le réflexe palpébral et le globe central.

\subsection{Nociception, douleur et conséquences}
\begin{align*}
N_{eff} &= \max_k\ N_k\,(1 - b_{terr(k)})\,(1 + \vp{sens_gain}\,S), \qquad
\nu^* = N_{eff}\,(1 - A_{méd})\,\Bigl(1 - \frac{\Delta^h}{\mathrm{MAC_{BAR}}^h + \Delta^h}\Bigr),\\
\dot\nu &= (\nu^* - \nu)/\tau_\nu,\qquad \dot\sigma = (\nu - \sigma)/\tau_{stress},\qquad
\dot S = k_S\,\nu\,(1 - S)(1 - E_{NMDA}) - S/\tau_S.
\end{align*}
$N_k$ : stimuli actifs (temps opératoires du scénario, stimulus test, douleur postopératoire décroissante) ;
$b$ : fraction bloquée par les blocs locorégionaux du territoire ; $A_{méd}$ : antinociception combinée.
Mouvement si $u_k < 1/(1 + (\Delta/N_{eff})^\gamma)$, $u_k$ tiré à chaque stimulus (graine explicite), patient inconscient.
Les conséquences passent par les effets JSON de \vp{anesthesia_model.json}. Une activation sympathique $\Delta s = 0{,}4\,\nu$ relève
le tonus de base \vp{s0} (FC) et, avec les coefficients mêmes du baroréflexe du moteur ($ees \times (1 + 0{,}4\,\Delta s)$,
$r_{sys}\times(1 + 0{,}5\,\Delta s)$, \vp{vol_shift} $+\ 0{,}2\cdot 85\cdot bw\cdot\Delta s$), l'inotropisme, les RVS et le recrutement veineux,
que le baroréflexe calcule relativement à \vp{s0}. $\nu$ retire du tonus vagal et réinitialise le baroréflexe vers le haut
(\vp{map_set} $+\ 20\,\nu$ mmHg, réaction de défense) ; sans cette réinitialisation, le baroréflexe opposerait une bradycardie à l'hypertension
de la stimulation. $\nu$ relève aussi la commande ventilatoire ; l'état de stress $\sigma$ augmente VO$_2$, VCO$_2$ et la glycémie. La sensibilisation $S$ amplifie les stimuli ultérieurs et la douleur postopératoire.
Indice de douleur du modèle $10\,\nu$ (échelle 0--10 non validée) quand le patient est conscient ; agitation au réveil si la
conscience est partielle et $\nu > 0{,}4$.

\subsection{Température, remplissage, hémorragie}
\[
\dot T = \frac{(T_0 - T_{amb})\,\frac{VO_2}{VO_{2,0}} - [1 + (g - 1)\,a\,v]\,(T - T_{amb})}{\tau} + \dot T_{réch},\qquad
\dot V_f = \text{débit} - V_f/\tau_{crist},
\]
$a = \mathrm{clip}(L/\vp{loc_mac}, 0, 1)$, $v = \mathrm{clip}\bigl((T - T_{vc})/(T_0 - T_{vc}), 0, 1\bigr)$ : l'excès de pertes s'éteint à l'approche de $T_{vc}$ = \vp{vasoconstriction_c} (vasoconstriction thermorégulatrice, plateau de Redondo 2012) ; $g$ = 3,5 et réchauffement 3,5 °C/h (recalés le 20/11/2026). Hypotension : le moteur stationnaire conclut au choc irréversible sous \vp{map_arrest} ; la session ne retient le décès qu'après \vp{shock_tol_min} (60 min) cumulées sous ce seuil. Le volume intravasculaire de cristalloïde s'ajoute à \vp{vol_shift}, les pertes sanguines
s'en retranchent (compensées par un remplissage transcapillaire vers 60\,\% de la perte, $\tau$ 2,9 h, une contraction splénique commandée par le sympathique, 6,8 mL/kg, et une vasoconstriction vasopressine--angiotensine qui monte au-delà de 25\,\% de perte ; 21/11/2026) ; $hb \leftarrow hb\cdot(VS - pertes)/(VS - pertes + V_f)$. VO$_2$ et VCO$_2$ suivent un Q$_{10}$ de 2 (effet JSON de $DT$).

\section{Modules et schémas}
\begin{tabularx}{\linewidth}{P{5.2cm}P{3.6cm}L}
\toprule
Fichiers & Schéma & Contenu \\
\midrule
\vp{modules/anesthesia/agents/*.json} & \vp{drug.schema.json} + bloc \vp{anesthesia} & 12 agents : dexmédétomidine, acépromazine, méthadone, fentanyl, propofol, alfaxalone, kétamine, midazolam, lidocaïne, bupivacaïne, atipamézole, naloxone. \\
\vp{modules/anesthesia/inhalants/*.json} & \vp{inhalant.schema.json} & isoflurane, sévoflurane. \\
\vp{modules/anesthesia/anesthesia_model.json} & \vp{anesthesia_model.schema.json} & constantes (chacune avec niveau de preuve et source), tissus, effets de la nociception et de la température. \\
\vp{modules/anesthesia/procedures/*.json} & \vp{procedure.schema.json} & ovariohystérectomie, orchiectomie (temps opératoires, territoires, douleur postopératoire). \\
\bottomrule
\end{tabularx}

Les modules d'anesthésie sont rangés à part : les catalogues existants (\vp{modules/drugs}) et le golden ne changent pas.
Les modules médicaments ordinaires (atropine, par exemple) peuvent aussi être chargés dans une session.

\section{Provenance des valeurs}
% Généré par tests/make_anesthesia_doc.py — ne pas éditer à la main.
\subsection{Agents IV, IM et locorégionaux}
\begin{tabularx}{\linewidth}{P{2.4cm}P{0.6cm}L}
\toprule
Agent & Niv. & Provenance (bloc \vp{anesthesia}) et notes \\
\midrule
Acépromazine & \evD & PK entièrement placeholder D (seule la t$\frac12$ IV de 7,1 h est publiée, Hashem 1992). Épargne de MAC plafonnée à 40 \% (halothane, Heard 1986, résumé : maximum dès 0,04 mg/kg) ; EC50 placeholder. Sédation placeholder. Blocage $\alpha$1 : veinodilatation (ven\_cap $\times$ (1 + 2$\cdot$E), borne du balayage) pour une PAM $-$12 \% chez le vigile (cible $-$17 \% sans changement de RVS, Rangel 2020 : écart signalé) ; hémoglobine $-$21 \% à effet maximal (CaO2 $-$21 \%, Rangel 2020 ; séquestration splénique, Sutil 2017). Pas d'analgésie. \\
Alfaxalone & \evC & PK : Vss 0,94 L/kg, Cl 26 mL/kg/min, Q 8,5 mL/kg/min (Dehuisser 2019, résumé) ; partage Vc/Vp placeholder D. ke0 placeholder D (classe). c\_imm calé pour que la dose d'induction publiée sans prémédication (2,6--2,8 mg/kg, Maney 2013, Hampton 2019) atteigne la perte de conscience avec la même marge que le propofol à 5,2 mg/kg (calage par session). Hémodynamique : PAM peu modifiée (Sarotti 2024, lu ; hypotension 1/14), FC +54 \% (Hampton 2019) : vasodilatation modérée placeholder D. Apnée : commande ventilatoire max(0, 1 $-$ 0,92$\cdot$EP) portée par le pic plasmatique, calée sur la dépendance à la vitesse d'injection (100 \% à 2 mg/kg/min contre 25 \% à 0,5 mg/kg/min, Bigby 2017) : hypopnée profonde en injection rapide après prémédication, ventilation conservée en injection lente. \\
Atipamézole & \evC & PK IM : Cl/F 554 mL/h/kg (9,23 mL/min/kg en supposant F = 1, D), Vz/F 1,49 L/kg (partage Vc/Vp D), ka 9,56 h$^{-1}$ calculé sur le Tmax de 0,35 h et la t$\frac12$ de 1,8 h (Cowan 2026, lu). Ki et ke0 placeholders D (éveil en 3 à 7 min après atipamézole IM, Vainio \& Vähä-Vahe 1990, résumé). Antagonisme compétitif de la dexmédétomidine (EC50 de toutes ses courbes multipliées par 1 + C/Ki). Rechute de sédation possible quand l'atipamézole s'élimine avant l'agoniste (Turunen 2019). \\
Bupivacaïne (bloc locorégional) & \evD & Module de bloc uniquement : délai, durée et efficacité placeholders D. PK systémique placeholder sans effet. À n'utiliser qu'en bloc locorégional (événement local\_block). \\
Dexmédétomidine & \evC & Clairance et Vss : Dent 2019 (résumé). Vc/Vp : partage placeholder D (Vc + Vp = Vss). IM : ka 5,96 h$^{-1}$ calculé sur le Tmax de 20,5 min (Aarnes 2023, résumé) avec t$\frac12$ 36 min ; F = 1 placeholder D. ke0 placeholder D (pic de sédation 12 min après IV, Alves 2024). Épargne de MAC (recalée le 20/11/2026) : EC50 3,0 ng/mL et pente 1,27 (max 0,8 imposé, D), ajustées sur les protocoles de Pascoe 2006 (charge puis CRI de 0,5 et 3 \textmu{}g/kg/h : $-$18 \% et $-$59 \%, cité par Acevedo-Arcique 2014, lu) aux concentrations que produit la PK du modèle à ces débits (1,1 et 6,7 ng/mL à 60 min). Les concentrations mesurées par Pascoe (0,198 et 1,903 ng/mL, résumé) sont 3 à 6 fois plus basses que celles du modèle : la PK (Dent 2019) n'est pas cohérente avec elles, écart non résolu (docs/VALIDATION\_HORS\_CALIBRAGE.tex). Hors échantillon : 2 \textmu{}g/kg + 2 \textmu{}g/kg/h, $-$43,4 \% (Acevedo-Arcique 2014, lu). Sédation : EC50 de l'ordre des concentrations au pic de sédation (Kuusela 2000 : 5,5 ng/mL à 10 \textmu{}g/kg) ; plafond 0,5 placeholder D (dose de propofol 5,2 $\to$ 2,9 mg/kg après 5 \textmu{}g/kg IM, Canfrán 2016). Antinociception : placeholder D. Effet biphasique : vasoconstriction périphérique ($\alpha$2B) portée par la concentration PLASMATIQUE (EP), sympatholyse centrale et tonus vagal par le site effet (E). Coefficients calés par petite grille (examples/anesthesia\_timing.rs) sur 10 \textmu{}g/kg IM à 15 min : modèle FC $-$58 \%, DC $-$57 \%, PAM +32 \% pour des cibles FC $-$55 \% et DC $-$59 \% (Congdon 2011) et PAM +36 \% (Alves 2024, 2 \textmu{}g/kg IV). Un bloc AV du 2e degré peut émerger du tonus vagal (av\_vagal du moteur) à plus forte dose, sans terme ad hoc ; il n'apparaît pas dans ce scénario de calage, alors que Karna 2022 le décrit chez tous les chiens (avec morphine). \\
Fentanyl & \evC & PK : Sano 2006 (résumé) : Cl 77,9 mL/min/kg, Vd 5 L/kg, t$\frac12$$\alpha$ 4,5 min, t$\frac12$$\beta$ 45,7 min ; Vc 4,05, Vp 0,95 L/kg et Q 116 mL/min/kg calculés (seule solution bicompartimentale compatible avec ces quatre valeurs). Épargne de MAC : plafond 84 \% et EC50 2,64 ng/mL (Machado 2022, résumé) ; pente 1 placeholder D. EC50 hémodynamique 3,65 ng/mL (FC, Machado 2022). Antinociception et sédation : EC50 reprises de la MAC, plafonds D ; ke0 D. Effets hémodynamiques et ventilatoires : forme et amplitudes du module opioide (classe, D pour le fentanyl), EC50 propre. \\
Kétamine & \evC & PK : Vc 0,371 L/kg, Vss 4,06 L/kg, Cl 58,2 mL/min/kg (Pypendop \& Ilkiw 2005, résumé) ; Q placeholder D. Épargne de MAC de l'isoflurane : Emax 92,2 \%, EC50 2,945 \textmu{}g/mL, $\gamma$ 3,01 (Pypendop 2007, résumé, sur les données de Solano 2006). c\_imm (perte de conscience) calé sur la dose d'induction de 6,1 mg/kg avec diazépam (Hampton 2019) : D. Antinociception, effet anti-hyperalgésique (NMDA) à faible concentration et ke0 : placeholders D. Anesthésie dissociative : réflexe palpébral conservé, globe central. Stimulation sympathique (s0 + 0,8$\cdot$E, plafond du balayage) : FC +23 \% après 6,1 mg/kg + midazolam pour une cible de +75 \% (Hampton 2019, kétamine-diazépam) ; le baroréflexe du moteur compense : écart signalé. Remarque : la contribution hypnotique (c\_imm) et l'épargne de MAC se cumulent dans la profondeur $\Delta$ : simplification signalée. \\
Lidocaïne & \evC & PK : Cl 27,5 mL/min/kg, Vss 1,38 L/kg (Ngo 1997, résumé) ; partage Vc/Vp et Q placeholders D. IM : ka 7,74 h$^{-1}$, fraction absorbée 91,9 \% (Wilcke 1983, résumé). Épargne de MAC : Emax ajusté sur Valverde 2004 (1,5 \textmu{}g/mL $-$18,7 \% ; 4,5 \textmu{}g/mL $-$43,3 \%, résumé) avec plafond 0,7 imposé (D) : EC50 3,15 mg/L, pente 1,36. Bloc locorégional (délai, durée, efficacité) et antinociception systémique : placeholders D. Toxicité non modélisée (le moteur n'a pas d'état convulsif) : extension tonique à 8,21 \textmu{}g/mL chez le chien (Wilcke 1983) ; ne pas dépasser cette concentration dans les scénarios. \\
Méthadone & \evC & SC : F 79 \% et ka 2,16 h$^{-1}$ calculé sur le Tmax de 1,26 h (Ingvast-Larsson 2010, résumé) ; IM non publiée : valeurs SC reprises (D). Épargne de MAC calée sur Credie 2010 (1 mg/kg : $-$48 \% à 2,5 h et $-$30 \% à 5 h ; 0,5 mg/kg : $-$35 \% et $-$15 \%, résumé, voie supposée IV). Antinociception : EC50 17 ng/mL (antinociception au-delà de 17 ng/mL, Amon 2021) ; plafond D. Sédation et ke0 : placeholders D. PK et PD hémodynamique reprises du module opioide 1.1.0 (méthadone, chien vigile). Sous halogéné, la dépression ventilatoire s'additionne à celle de l'agent inhalé (pas de terme de synergie). \\
Midazolam & \evC & PK : Vd 3,0 L/kg, Cl 27 mL/kg/min (Court \& Greenblatt 1992, résumé) ; partage Vc/Vp, Q, IM et ke0 placeholders D. Sédation plafonnée à 0,5 : la cible de propofol à l'intubation passe de 3,0 à 1,5 \textmu{}g/mL avec 0,2 mg/kg (Minghella 2016, résumé). Épargne de MAC : enflurane $-$55 \% à 949 ng/mL (Hall 1988, résumé) ; plafond 0,6 et EC50 D. Peu d'effets hémodynamiques ; légère dépression ventilatoire placeholder D. L'excitation paradoxale du chien vigile n'est pas modélisée. \\
Naloxone & \evC & t$\frac12$ IV 37 min et Cmax 18,8 ng/mL après 0,04 mg/kg (Wahler 2019, résumé) : Vc $\approx$ 2,1 L/kg et Cl $\approx$ 58 mL/min/kg calculés (approximation, C/D) ; Vp, IM, Ki et ke0 placeholders D (antagonisation de la méthadone en < 5 min après 0,1 mg/kg IM, Mills 2025, lu). Renarcotisation possible si la naloxone s'élimine avant l'opioïde (t$\frac12$ 37 min contre 3,9 h pour la méthadone). \\
Propofol & \evC & Cl 58,6 mL/kg/min et Vss 4,9 L/kg (Nolan 1993, résumé) ; partage Vc 0,8 / Vp 4,1 L/kg et Q placeholders D (Vc non publié). ke0 0,696 min$^{-1}$ (Lee 2009, résumé). c\_imm 6,2 mg/L : ajusté sur la réduction de MAC du sévoflurane (1,3 / 2,5 / 4,4 \textmu{}g/mL $\to$ $-$20,5 / $-$43 / $-$68 \%, Singsank-Coats 2015, résumé ; contribution Ce/c\_imm $\approx$ 0,16 MAC par \textmu{}g/mL) ; il place la perte de conscience (indice 0,4) vers 2,5 \textmu{}g/mL, cohérent avec Musk 2005 (65 \% d'inductions à 2,5 \textmu{}g/mL, 100 \% à 3,5) et la cible de 3 \textmu{}g/mL de Beths 2001. Vasodilatation (r\_sys) et émoussement du baroréflexe (baro\_gain, mécanisme retenu, forme D) calés ensemble pour PAM $-$31 \% au nadir et FC +62 \% (88 $\to$ 61 mmHg, 65 $\to$ 105/min, Cattai 2018, lu, après acépromazine + méthadone) ; la tachycardie vient du baroréflexe du moteur. Effets hémodynamiques portés par la concentration PLASMATIQUE (EP) : nadir de PAM à 55 s, avant l'équilibre du site effet (intubation à 140 s, Cattai 2018). Apnée d'induction : commande ventilatoire max(0, 1 $-$ 1,23$\cdot$EP), portée par le pic PLASMATIQUE : une injection rapide l'abolit, et le CO2 ne la relance qu'une fois la concentration redescendue. Calée sur la dépendance à la vitesse (apnée chez 100 \% des chiens à 4 mg/kg/min contre 25 \% à 1 mg/kg/min, Bigby 2017 ; 59 à 64 \% à 4 mg/kg, Murison 2001) : 5,2 mg/kg en 30 s sans prémédication donnent 56 s d'apnée, 1,8 mg/kg à 1 mg/kg/min après prémédication aucune. La vitesse d'injection se règle par la durée du bolus (30 s par défaut, iv\_push\_s du modèle). \\
\bottomrule
\end{tabularx}

\subsection{Halogénés}
\begin{tabularx}{\linewidth}{P{2.4cm}P{0.6cm}L}
\toprule
Agent & Niv. & Valeurs et provenance \\
\midrule
Isoflurane & \evC & MAC 1.3 \%, sang/gaz 1.4, tissu/sang 1.57 / 2.9 / 45.0. MAC 1,3 \% (1,28 à 1,39 \% : Steffey \& Howland 1977, Kazama 1988, Hofmeister 2008, résumés). Sang/gaz canin 1,40 (Soares 2012, lu). Cerveau/sang 1,57 (humain, Yasuda 1989) ; muscle/sang 2,9 et graisse/sang 45 : ordres de grandeur humains classiques, NON LUS (D). Effets en fonction de M = pression artérielle / MAC : Ees VG et VD $-$23,5 \% par MAC (Hoka 1995 : $-$23 \% à 1 MAC, $-$48 \% à 2 MAC) ; vasodilatation r\_sys = max(0,1 ; 1 $-$ c$\cdot$M\textasciicircum{}p) et émoussement du baroréflexe baro\_gain $\times$ (1 $-$ b$\cdot$M) calés par session sur le sévoflurane de Bernard 1990 (PAM $-$22 \% à 1,2 MAC, $-$38 \% à 2 MAC, DC $-$17 \% à 2 MAC ; isoflurane « similaire »). La dépression du baroréflexe par les halogénés est le mécanisme retenu (non chiffrée dans les sources lues : forme D) ; commande ventilatoire max(0, 1 $-$ a$\cdot$M\textasciicircum{}n) calée pour PaCO2 46 mmHg à 1 MAC en ventilation spontanée (Steffey 1993) et apnée à 2,6 MAC (Galloway 2004). Baroréflexe non déprimé par l'agent (placeholder : effet non sourcé). \\
Sévoflurane & \evC & MAC 2.2 \%, sang/gaz 0.664, tissu/sang 1.7 / 3.1 / 48.0. MAC 2,2 \% (2,10 à 2,36 \% : Kazama 1988, Hofmeister 2008, Yamashita 2012, Murahata 2018). Sang/gaz canin 0,664 (Soares 2012, lu). Cerveau/sang 1,70 (humain, Yasuda 1989) ; muscle/sang 3,1 et graisse/sang 48 : ordres de grandeur humains NON LUS (D). Ees $-$20 \% par MAC (Hoka 1995 : $-$16 \% à 1 MAC, $-$41 \% à 2 MAC) ; vasodilatation max(0,1 ; 1 $-$ c$\cdot$M\textasciicircum{}p) et émoussement du baroréflexe (1 $-$ b$\cdot$M, mécanisme retenu, forme D) calés sur Bernard 1990 (PAM $-$22 \% à 1,2 MAC, $-$38 \% à 2 MAC, DC $-$17 \% à 2 MAC) ; commande ventilatoire calée pour l'apnée à 3,45 MAC (Galloway 2004), avec la même PaCO2 à 1 MAC que l'isoflurane (hypothèse D, aucune valeur canine publiée trouvée). \\
\bottomrule
\end{tabularx}

\subsection{Constantes du modèle}
\begin{tabularx}{\linewidth}{P{3.6cm}P{1.2cm}P{0.6cm}L}
\toprule
Constante & Valeur & Niv. & Source \\
\midrule
\vp{co2_store_ml_mmhg_kg} & 1.0 & \evD & Placeholder : capacité de stockage du CO2 choisie pour une montée de PaCO2 de l'ordre de 4 à 5 mmHg/min en apnée (VCO2 4,6 mL/kg/min du patient). À sourcer sur une étude d'apnée chez le chien. Cohérence vérifiée : en apnée totale sous O2, PaCO2 de 96 mmHg à l'arrêt (DeBehnke 1995 : 93 $\pm$ 19 mmHg à 11,4 min, résumé). \\
\vp{frc_ml_kg} & 45.0 & \evC & CRF 45 $\pm$ 8 mL/kg chez le chien sous halothane, curarisé, ventilé (Larsson 1992, Acta Anaesthesiol Scand 36:146, résumé) ; 30 à 75 mL/kg selon la posture et la sédation chez le chien vigile (Rozanski 2010, PMC2801309, lu). \\
\vp{atelectasis_aa_mmhg} & 110.0 & \evC & Gradient A-a ajouté sous anesthésie, proportionnel à la profondeur et à (FiO2 $-$ 0,21)/0,79 : effet du shunt d'atélectasie, négligeable à l'air. Calé le 20/11/2026 : PaO2 547 [495--616] mmHg sous isoflurane 1,3--1,5 \% en ventilation contrôlée et O2 (Itami 2019, doi:10.1371/journal.pone.0213414, lu ; modèle 548) ; 523 (513--539) mmHg sous sévoflurane 1,3 MAC en ventilation spontanée (Kato 2024, doi:10.5455/OVJ.2024.v14.i5.20, lu ; modèle 523) ; 570--579 chez Mosing 2025 (lu). Avant : 630 à 660 mmHg. \\
\vp{vd_ml_kg} & 3.0 & \evD & Placeholder : espace mort anatomique + appareil en ventilation contrôlée. \\
\vp{circuit_l} & 3.0 & \evD & Placeholder : volume d'un circuit filtre adulte. \\
\vp{etco2_gap_mmhg} & 3.0 & \evD & Placeholder : écart PaCO2 $-$ EtCO2 chez le sujet sain. \\
\vp{flowby_fio2} & 0.4 & \evD & Placeholder : FiO2 obtenue en flow-by. \\
\vp{loc_mac} & 0.65 & \evC & MAC-awake du chien : rapports 0,77 (isoflurane 1,0/1,3 \%) et 0,62 (sévoflurane 1,3/2,1 \%), Hofmeister 2008 (résumé) ; moyenne arrondie. \\
\vp{awareness_hill} & 3.0 & \evC & Pente de la conscience sous la MAC-awake, awareness = 1 $-$ (indice/loc\_mac)\textasciicircum{}n. Calée le 20/11/2026 sur les délais de réveil (décubitus sternal) : isoflurane FE 1,3 \% ($\approx$ 1,0 MAC) 2 h, sternal 14,7--18,0 min (Rastas 2025, PMC11781748, lu ; modèle 17,2) ; sévoflurane $\approx$ 105--120 min, extubation 7 (5--7) et sternal 9,5 (7,25--11,75) min (Lozano 2009, résumé ; modèle à 1,45 MAC : 9,0 et 12,1). Avec n = 1 (avant) : sternal à 104 min. \\
\vp{swallow_mac} & 0.5 & \evC & MAC-extubation du sévoflurane : rapport 0,51 $\pm$ 0,08 (le chien lève la tête ou mâchonne la sonde), Murahata 2018 (PMC6068302, lu). \\
\vp{palpebral_mac} & 0.9 & \evD & Placeholder : perte du réflexe palpébral. \\
\vp{jaw_mac} & 1.2 & \evD & Placeholder : perte du tonus de la mâchoire. \\
\vp{surgical_mac} & 1.0 & \evD & Placeholder : début du plan chirurgical (1 MAC : 50 \% d'immobilité au stimulus supramaximal). \\
\vp{deep_mac} & 1.8 & \evD & Placeholder : plan profond (globe recentré). \\
\vp{excessive_mac} & 2.5 & \evD & Placeholder : surdosage (mydriase). \\
\vp{mac_bar} & 1.6 & \evC & MAC-BAR du sévoflurane : 3,33 \% pour une MAC de 2,10 \% (rapport 1,59, Yamashita 2012, résumé) ; rapport 2,05 chez Murahata 2018. \\
\vp{bar_hill} & 3.0 & \evD & Placeholder : pente de la suppression de la réponse autonome. \\
\vp{move_hill} & 10.0 & \evD & Placeholder : pente de la courbe d'immobilité (relation dose-réponse raide de la MAC). \\
\vp{tau_nu_s} & 15.0 & \evD & Placeholder : constante de temps de la réponse nociceptive. \\
\vp{tau_stress_min} & 20.0 & \evD & Placeholder : constante de temps de la réponse de stress (VO2, glycémie). \\
\vp{sens_rate_per_min} & 0.05 & \evD & Placeholder : vitesse de sensibilisation centrale sous nociception non contrôlée. \\
\vp{sens_tau_h} & 24.0 & \evD & Placeholder : retour de la sensibilisation. \\
\vp{sens_gain} & 0.5 & \evD & Placeholder : amplification maximale de la nociception par la sensibilisation (+50 \%). \\
\vp{t0_c} & 38.5 & \evD & Placeholder : température de référence du chien. \\
\vp{t_amb_c} & 21.0 & \evD & Placeholder : température du bloc opératoire. \\
\vp{tau_cool_h} & 12.0 & \evD & Placeholder : constante thermique ; avec loss\_gain\_anesth = 2 et VO2 $-$20 \%, refroidissement initial $\approx$ 1,75 °C/h. \\
\vp{loss_gain_anesth} & 3.5 & \evC & Pertes thermiques multipliées sous anesthésie (vasodilatation, perte du comportement), au-dessus du seuil de vasoconstriction. Calé le 20/11/2026 avec vasoconstriction\_c sur Redondo 2012 (1525 chiens, résumé) : $-$2,0, $-$2,6 et $-$2,9 °C à 60, 120 et 180 min (modèle $-$2,0, $-$2,6, $-$2,8 ; isoflurane 1,3 MAC, salle 21 °C) ; chute rapide puis plateau après 60 min (Sage 2023, 1097 chiens, résumé). Avant : 2,0 sans plateau ($-$4 °C à 3 h). \\
\vp{vasoconstriction_c} & 35.5 & \evC & Température où l'excès de pertes de l'anesthésie s'éteint (vasoconstriction thermorégulatrice) : décroissance linéaire de l'excès entre t0\_c et ce seuil. Calé le 20/11/2026 avec loss\_gain\_anesth sur le plateau de Redondo 2012 (35,8 $\pm$ 1,5 °C à 180 min, résumé) et Sage 2023 (résumé). Forme D. \\
\vp{warming_c_per_h} & 3.5 & \evC & Apport d'un réchauffement actif (matelas à eau + air pulsé). Calé le 20/11/2026 : température finale 38,2 $\pm$ 1,0 et 38,7 $\pm$ 0,5 °C avec réchauffement contre 35,4--35,7 sans, après $\approx$ 2 h d'isoflurane (Rastas 2025, PMC11781748, lu ; modèle 38,4 à 3 h). Avant : 1,5 °C/h (insuffisant avec les pertes recalées). \\
\vp{mac_per_c} & 0.05 & \evD & Placeholder : réduction de MAC par °C d'hypothermie (ordre de grandeur humain non lu). \\
\vp{blood_volume_ml_kg} & 85.0 & \evD & Volume sanguin de 85 mL/kg, déjà utilisé par le baroréflexe du moteur (MODEL\_SPEC \S{}3). \\
\vp{crystalloid_tau_min} & 30.0 & \evD & Placeholder : sortie vasculaire des cristalloïdes. \\
\vp{refill_frac} & 0.6 & \evC & Fraction du volume perdu restituée au plasma par remplissage transcapillaire ; avec refill\_tau\_h, 30 \% à 2 h (Joles 1989, chiens conscients, résumé) et 56 \% à 8 h (Miller \& Dale 1978, résumé). \\
\vp{refill_tau_h} & 2.9 & \evC & Constante du remplissage transcapillaire, ajustée avec refill\_frac sur Joles 1989 et Miller \& Dale 1978 (résumés). \\
\vp{spleen_ml_kg} & 6.8 & \evC & Sang mobilisable par la rate : 8 \% du volume sanguin en hémorragie chez le chien conscient (Guntheroth \& Mullins 1963, résumé) ; 2,4--3,2 mL/kg dans les réflexes et hémorragies modérées (Shoukas 1981, Carneiro \& Donald 1977). \\
\vp{spleen_tau_min} & 1.0 & \evD & Placeholder : contraction splénique en une à quelques minutes. \\
\vp{spleen_s_span} & 0.3 & \evD & Placeholder : hausse du tonus sympathique qui vide entièrement la rate. \\
\vp{spleen_hct_ratio} & 1.8 & \evD & Placeholder : sang splénique concentré (Ht 28,9 \% chez les splénectomisés contre 39,8 \% chez les témoins après 50 \% de saignée, Hoekstra 1988, résumé). \\
\vp{pressor_gain} & 0.6 & \evC & Hausse maximale des RVS par vasopressine et angiotensine II en hémorragie. Calée le 21/11/2026 sur le chien conscient : PAM $-$10 mmHg à 17 mL/kg (Shen 1990 : $-$15 $\pm$ 3 à 20 mL/kg), $-$23 mmHg à 30 mL/kg (Thrasher 1998 : $-$35 $\pm$ 10), DC $-$41 \% à 25 \% (Francis 2004 : $-$40 \%, résumés). Le blocage de l'angiotensine aggrave la chute (losartan : $-$40 \% au lieu de $-$25 \%, Francis 2004 ; saralasine : PAM 102 $\to$ 80, Spielman 1975). \\
\vp{pressor_lost_frac} & 0.25 & \evC & Mi-effet presseur vers 25 \% du volume sanguin : vasopressine 6, 200 puis 991 pg/mL à 10, 20 et 30 mL/kg (Wang 1988, résumé). \\
\vp{pressor_hill} & 4.0 & \evD & Placeholder : montée raide de la vasopressine entre 10 et 30 mL/kg. \\
\vp{iv_push_s} & 30.0 & \evC & Durée d'injection par défaut d'un bolus IV : 30 s, comme l'induction au propofol de Cattai 2018 (PMC6134702, lu). La vitesse d'injection change l'incidence de l'apnée (Murison 2001, Bigby 2017) et l'hypotension : elle est réglable à chaque injection. \\
\vp{o2_debt_tol_ml_kg} & 20.0 & \evD & Placeholder : dette d'O2 tolérée avant l'arrêt hypoxique (réserves, anaérobiose), environ 4 min de VO2 totale (5,5 mL/kg/min) : l'arrêt survient alors quelques minutes après le début d'une apnée à l'air ambiant. Le moteur stationnaire, lui, conclut au décès dès que VO2 dépasse DO2. Délai d'arrêt du modèle en apnée totale sous O2 : 10,3 min depuis l'atélectasie du 20/11/2026 (11,6 avant ; DeBehnke 1995 : 11,4 $\pm$ 2,4 min, chiens anesthésiés, sonde clampée). \\
\vp{o2_debt_tau_min} & 10.0 & \evD & Placeholder : constante de remboursement de la dette d'O2 une fois l'apport rétabli. \\
\vp{shock_tol_min} & 60.0 & \evC & Durée cumulée sous map\_arrest (PAM < 40 mmHg) avant le choc irréversible : 1 à 3 h chez le chien maintenu vers 40 mmHg (Morales 1999, Granot 1985, résumés, déjà cités pour map\_arrest), borne basse retenue. Le moteur stationnaire conclut au décès dès PAM < 40 ; la session tolère cette hypotension pendant ce délai. Chiens survivant à PAM 40--44 mmHg et DC 0,6 L/min sous isoflurane (Cardillo 2024, PMC11551537, lu). Ajouté le 20/11/2026 (docs/VALIDATION\_HORS\_CALIBRAGE.tex). \\
\vp{do2_crit_ml_kg_min} & 9.8 & \evC & Apport d'O2 critique 9,8 mL/kg/min chez le chien anesthésié, hypoxie hypoxique et anémique (Cain 1977, J Appl Physiol 42:228, résumé) ; 7,7 mL/kg/min en hémorragie (Schertel 1994, résumé). \\
\vp{blood_o2_store_frac} & 0.5 & \evD & Fraction du sang dont le contenu en O2 suit l'alvéole pendant l'apnée (réserve sanguine, pente de la courbe de dissociation) : placeholder calé sur le délai d'arrêt par asphyxie (modèle 11,6 min sous O2, DeBehnke 1995 : 11,4 $\pm$ 2,4 min, résumé) ; sans effet à l'équilibre. \\
\bottomrule
\end{tabularx}

\subsection{Scénarios chirurgicaux}
\textbf{Orchiectomie (chien)} (\evD) : Incision prescrotale (0.8, 60 s) ; Ligature du premier cordon spermatique (1.1, 120 s) ; Ligature du second cordon spermatique (1.1, 120 s) ; Fermeture (0.5, 300 s). Douleur postopératoire 0.4, constante 8 h. Intensités relatives : placeholders pédagogiques D. Le bloc intratesticulaire se pose sur le territoire « spermatic\_cord ».

\textbf{Ovariohystérectomie (chienne)} (\evD) : Incision cutanée et de la ligne blanche (0.8, 60 s) ; Extériorisation de la corne utérine (0.7, 120 s) ; Rupture du ligament suspenseur et ligature du pédicule ovarien (1.2, 180 s) ; Second pédicule ovarien (1.2, 180 s) ; Ligature du col utérin (0.9, 120 s) ; Fermeture de la paroi et de la peau (0.5, 600 s). Douleur postopératoire 0.6, constante 12 h. Intensités relatives (1 = stimulus supramaximal de la définition de la MAC) : placeholders pédagogiques D qui ordonnent les temps opératoires (le pédicule ovarien est le temps le plus stimulant, enseignement classique). Aucune valeur publiée.



\section{Calibration}
Le banc \vp{crates/vetsim-api/examples/anesthesia_calibrate.rs} rejoue chaque cible publiée dans le modèle complet
(session, baroréflexe et décompensation compris). Avec \vp{--fit}, il cherche les coefficients par balayage déterministe et
les imprime en JSON ; ils sont reportés à la main dans le fichier de coefficients ci-dessous, jamais modifiés automatiquement dans \vp{modules/}.
Coefficients calés (\vp{tests/anesthesia_modules/coefficients.json}, générateurs \vp{gen_agents.py} et \vp{gen_inhalants.py}) :
\begin{itemize}
\item halogénés : vasodilatation $\max(0{,}1;\,1 - c\,M^p)$ et émoussement du baroréflexe $1 - b\,M$ sur trois cibles de Bernard (1990) ;
commande ventilatoire $\max(0;\,1 - a\,M^n)$ sur la PaCO$_2$ de 46,2 mmHg à 1 MAC (Steffey 1993) et l'apnée à l'indice publié (Galloway 2004) ;
\item dexmédétomidine : petite grille bornée (\vp{examples/anesthesia_timing.rs}) sur vasoconstriction plasmatique, sympatholyse et tonus vagal ;
la bissection automatique (\vp{--dex}) explore des résistances extrêmes très coûteuses ;
\item acépromazine : veinodilatation ; propofol : vasodilatation et émoussement du baroréflexe portés par la concentration plasmatique
(grille, la relation n'est pas monotone) ;
\item $c_{imm}$ de l'alfaxalone et de la kétamine (même marge d'induction que le propofol à 5,2 mg/kg), stimulation sympathique de la kétamine,
épargne de MAC de la méthadone (ajustement direct sur $C_e$, qui ne dépend que de la PK).
\end{itemize}
Tous les bolus IV de la calibration sont injectés en \SI{30}{s} (\vp{iv_push_s}).

Sortie du banc sur les modules livrés (chien sain de \SI{20}{kg}) :
\VerbatimInput[fontsize=\scriptsize,breaklines=true,breakanywhere=true]{anesthesie_calibration.txt}

\paragraph{Écarts résiduels.}
\begin{itemize}
\item Halogénés à 2 MAC : PAM $-48$\,\% contre $-38$\,\% et pas de tachycardie (+1\,\% contre +54\,\%) ; le moteur ne peut pas associer,
avec un baroréflexe émoussé, une FC élevée et un débit en baisse ; le compromis retenu privilégie la PAM à 1,2 MAC et le débit à 2 MAC.
\item Dexmédétomidine : FC $-37$\,\% contre $-55$\,\% ; le bloc AV du 2\textsuperscript{e} degré, qui émerge du tonus vagal comme chez Karna (2022),
plafonne la fréquence ventriculaire du moteur.
\item Kétamine-midazolam : FC +23\,\% contre +75\,\% ; le baroréflexe compense la stimulation sympathique (plafond du balayage).
\item Méthadone 0,5 mg/kg : épargne de 30\,\% à 2,5 h contre 35\,\%.
\end{itemize}


\section{Validation}
Tests \vp{crates/vetsim-core/tests/anesthesia.rs} (21) et \vp{crates/vetsim-api/tests/anesthesia.rs} (1). Chaque test compare un
scénario à un témoin sans le facteur étudié, pour pouvoir échouer pour la raison qu'il vérifie.
\begin{tabularx}{\linewidth}{P{6.2cm}L}
\toprule
Test & Propriété \\
\midrule
\vp{session_without_drugs_reproduces_eval} & invariant d'équilibre ($10^{-9}$). \\
\vp{apnea_raises_paco2_and_preoxygenation_delays_desaturation} & PaCO$_2$ monte en apnée ; délai de désaturation plus de 3 fois plus long après 3 min d'O$_2$. \\
\vp{sevoflurane_equilibrates_faster_than_isoflurane} & $F_A/F_I$ plus haut à 5 min (solubilité). \\
\vp{low_fresh_gas_flow_slows_deepening} & bas débit de gaz frais : approfondissement plus lent. \\
\vp{isoflurane_depresses_blood_pressure_and_ventilation_dose_dependently} & PAM, PaCO$_2$, $\Delta$ dépendants de la dose. \\
\vp{fentanyl_spares_mac} & épargne de MAC, $\Delta$ augmenté à vaporisateur égal. \\
\vp{surgical_stimulus_activates_sympathetic_and_analgesia_blunts_it} & le pédicule ovarien élève FC et PAM ; fentanyl et bloc atténuent $\nu$. \\
\vp{light_plane_moves_under_supramaximal_stimulus} & mouvement sous 1 MAC, pas en plan profond. \\
\vp{fast_propofol_injection_causes_apnea_slow_does_not} & apnée en injection rapide, pas ou moins en injection lente. \\
\vp{breathing_resumes_when_co2_rises_after_hyperventilation} & apnée au sevrage d'une hyperventilation, reprise quand la PaCO$_2$ remonte. \\
\vp{untreated_apnea_desaturates_then_causes_hypoxic_arrest} & apnée totale : désaturation puis arrêt hypoxique. \\
\vp{sevoflurane_recovers_faster_than_isoflurane} & retour de la déglutition plus précoce (ordre de Lopez 2009). \\
\vp{atipamezole_reverses_dexmedetomidine} & sédation levée, FC remontée. \\
\vp{premedication_reduces_propofol_requirement} & 2,6 mg/kg de propofol : inconscience seulement après prémédication. \\
\vp{uncontrolled_nociception_sensitizes_and_ketamine_limits_it} & sensibilisation, freinée par la kétamine. \\
\vp{anesthesia_cools_and_warming_limits_it} & hypothermie peropératoire, limitée par le réchauffement. \\
\vp{hemorrhage_lowers_pressure_and_hematocrit_dilutes_with_fluids} & hémorragie, remplissage, hémodilution. \\
\vp{deterministic_for_a_seed}, \vp{record_and_milestones}, \vp{invalid_actions_are_rejected}, \vp{modules_load} & déterminisme, feuille et jalons, erreurs (dont l'intubation d'un chien vigile), validation. \\
\bottomrule
\end{tabularx}

\section{API et interface}
Messages JSON (\vp{vetsim-api}, donc UniFFI et WebAssembly) :
\begin{Verbatim}
{"op": "anesthesia_list"}
{"op": "anesthesia_start",  "patient": "p1", "options": {"seed": 1, "procedure": "ovariohysterectomy"}}
{"op": "anesthesia_event",  "options": {"session": "a2", "event": {"type": "bolus", "agent": "propofol",
                                        "mg_per_kg": 4, "route": "iv"}}}
{"op": "anesthesia_step",   "options": {"session": "a2", "dt_s": 5, "n": 12, "all": true, "waves_s": 4}}
{"op": "anesthesia_step",   "options": {"session": "a2", "waves_s": 6, "seed": 1, "pcg": {"site": "apex"}}}
{"op": "anesthesia_record", "options": {"session": "a2"}}
\end{Verbatim}
Événements : \vp{bolus} (\vp{duration_s} optionnelle), \vp{cri}, \vp{vaporizer}, \vp{fresh_gas}, \vp{circuit}, \vp{intubate}, \vp{extubate}, \vp{ventilation},
\vp{oxygen}, \vp{stimulus}, \vp{procedure_step}, \vp{end_surgery}, \vp{local_block}, \vp{fluids}, \vp{fluid_bolus}, \vp{hemorrhage}, \vp{warming}.
\vp{loadModules} accepte un bloc \vp{anesthesia} (\vp{model}, \vp{agents}, \vp{inhalants}, \vp{procedures}).

L'onglet « Anesthésie » de l'application web (\vp{web/anesthesia.js}) affiche un moniteur (ECG D~II et pression artérielle simulés
par le moteur, capnogramme schématique, valeurs, alarmes), les signes cliniques, trois tendances, les actions et la feuille d'anesthésie ;
vitesse ×1, ×10 ou ×60.

\textbf{Auscultation en temps réel} (20/11/2026). Avec l'option \vp{pcg} (\vp{true}, ou un objet avec \vp{site} : \vp{base} ou \vp{apex}, et \vp{fs}),
\vp{anesthesia_step} renvoie avec les courbes un son cardiaque (\vp{pcg.wav_base64}, \vp{fs}, instants des battements, grade et type de souffle)
synthétisé par le module auscultation (\vp{pcg::patient_pcg}) sur l'évaluation courante de la session, avec la même graine et la même durée
que l'ECG : mêmes battements (test d'API \vp{anesthesia_step_streams_heart_sounds...} : chaque battement du son a une onde R à moins
de 0,12 s ; avec une autre graine, le test échoue). Sans configuration \vp{auscultation} chargée, la demande est refusée. Dans l'onglet,
le bouton « Stéthoscope » (site apex gauche ou base) découpe le son aux mêmes ondes R que l'ECG et le joue quand le segment passe sous la
tête de lecture du moniteur (temps réel, quelle que soit la vitesse de simulation, fondus de 5 ms aux raccords) ; il entend donc les effets
des agents, de la profondeur, de la volémie et de la maladie du patient tels que le moteur les calcule.

\section{Limites connues}
\begin{itemize}
\item Baroréflexe en régime permanent à chaque pas (pas de retard propre) : acceptable pour des pas de quelques secondes.
\item Le couplage d'ischémie coronaire du moteur peut basculer en « spirale de décompensation » (décès \vp{vicious_circle}) quand une
vasodilatation forte et un baroréflexe émoussé se cumulent (surdosage de propofol injecté très vite, par exemple) ; comportement du moteur,
non modifié ici, à garder en tête en enseignement.
\item Émoussement du baroréflexe par les halogénés et le propofol : mécanisme retenu pour reproduire les cibles publiées, sans valeur canine
mesurée de la sensibilité baroréflexe (forme \evD).
\item Aucune synergie spécifique opioïde--halogéné sur la ventilation (effets multiplicatifs simples) ; la PaCO$_2$ de 55 mmHg
après morphine sous isoflurane (Steffey 1993) n'est pas une cible du modèle.
\item Contribution hypnotique et épargne de MAC de la kétamine cumulées dans $\Delta$.
\item Propofol : $c_{imm}$ calé sur l'épargne de MAC place la perte de conscience vers \SI{4}{mg/L}, au-dessus des cibles d'induction
publiées (2,5--3,5 µg/mL prédits) : écart signalé, non corrigé.
\item Toxicité de la lidocaïne (convulsions) non représentée ; ESV de stress non branchées ; pas de modèle de voies aériennes
(obstruction, intubation sélective), ni de chat.
\item Intensités des temps opératoires, constantes de nociception, de température et de stockage des gaz : placeholders \evD.
\end{itemize}

\end{document}
